\documentclass[11pt]{article}
\newcommand{\documenttype}{Cover letter}
\usepackage[letterpaper,margin=1in]{geometry}
\usepackage[T1]{fontenc}
\usepackage[utf8]{inputenc}
\usepackage{newtxtext,newtxmath}
\usepackage{microtype}
\usepackage{xcolor}
\usepackage{enumitem}
\usepackage{tabularx}
\usepackage{titlesec}
\usepackage{fancyhdr}
\usepackage{hyperref}
\definecolor{accent}{HTML}{315B8A}
\definecolor{link}{HTML}{5B4F7A}
\hypersetup{colorlinks=true,urlcolor=link,linkcolor=link,citecolor=link,
  pdfauthor={Phurailatpam Hemantakumar}}
\setlength{\parindent}{0pt}
\setlength{\parskip}{6pt}
\setlength{\headheight}{14pt}
\setlength{\emergencystretch}{1.5em}
\titleformat{\section}{\large\bfseries\color{accent}}{}{0pt}{}
\titlespacing*{\section}{0pt}{11pt}{5pt}
\titleformat{\subsection}{\normalsize\bfseries}{}{0pt}{}
\titlespacing*{\subsection}{0pt}{8pt}{3pt}
\setlist[itemize]{leftmargin=15pt,itemsep=3pt,topsep=3pt,parsep=0pt}
\setlist[enumerate]{leftmargin=19pt,itemsep=7pt,topsep=4pt,parsep=0pt}
\pagestyle{fancy}
\fancyhf{}
\fancyhead[L]{\small Phurailatpam Hemantakumar}
\fancyhead[R]{\small\documenttype}
\fancyfoot[R]{\small\thepage}
\renewcommand{\headrulewidth}{0.3pt}
\renewcommand{\footrulewidth}{0pt}
\newcommand{\firstheading}[2]{%
  \thispagestyle{plain}%
  {\LARGE\bfseries\color{accent} #1}\par
  \vspace{3pt}{#2}\par
  \vspace{4pt}{\color{accent}\hrule}\vspace{8pt}}
\newcommand{\dated}[2]{\par\noindent\begin{tabularx}{\linewidth}{@{}X r@{}}#1 & #2\end{tabularx}\par}

\hypersetup{pdftitle={Cover letter | Phurailatpam Hemantakumar | CCA}}
\begin{document}
\firstheading{Phurailatpam Hemantakumar}{%
\href{mailto:hemantaph@link.cuhk.edu.hk}{hemantaph@link.cuhk.edu.hk}\quad
\href{https://hemantaph.com/}{hemantaph.com}}

31 October 2026

Flatiron Research Fellow Search Committee\\
Center for Computational Astrophysics\\
Flatiron Institute, Simons Foundation

Dear Members of the Search Committee,

I am applying for the Flatiron Research Fellow position at the Center for Computational Astrophysics. I completed my PhD in Physics at The Chinese University of Hong Kong in 2026. My PhD research covered two main areas: gravitational-wave lensing statistics, including signal-to-noise ratio (SNR) and detection-probability calculations, and eccentric waveform modelling and data analysis.

I lead the development of \texttt{ler}, an open-source framework for simulating lensed and unlensed gravitational-wave populations and predicting their detectable rates. It combines source and lens populations, strong-lensing selection, semi-analytic lens-equation solutions, and detector selection. Its code and documentation are maintained with contributions and feedback from users.

I developed \texttt{gwsnr} for efficient SNR calculations using partial-scaling interpolation and parallelised waveform inner products, as well as rapid detection-probability estimation using neural networks. Both packages have supported collaborative publications.

At CCA, my lensing-statistics programme would combine star-formation-informed merger-rate histories (metallicity-weighted star formation convolved with a formation-to-merger delay-time distribution), updated lens populations, and tested detection probabilities. I would quantify how uncertainties in the source and lens populations affect predicted lensing rates and population-based priors used in Bayesian lensing hypothesis testing. In a separate proposed project, I would examine how these population uncertainties affect cosmological inference from the lensed-event fraction and time-delay distribution. I am particularly interested in collaboration with Will Farr and Jack Heinzel on population inference and selection effects, and in exchanges with CCA's cosmology researchers.

In my waveform-modelling and data-analysis project, I developed the low-eccentricity \texttt{TaylorF2Ecck} inspiral waveform with collaborators, implemented it in \texttt{LALSuite}, and applied it to Bayesian analyses of GW170817 and GW190425. I plan to improve the waveform models and investigate what the resulting constraints imply for compact-binary formation. I am interested in collaboration with Max Isi on waveform systematics and inference from individual events.

In related work on individual events, I investigated the degeneracy between eccentricity and line-of-sight acceleration in GW190814. Separately, building on our earlier joint detectability forecasts, I am developing Bayesian posterior-odds tests for strongly lensed gamma-ray bursts, adapting the gravitational-wave lensing framework of Hannuksela et al. (2026) to GRB sky-localisation and time-delay measurements, with or without detectable gravitational-wave counterparts.

I have mentored undergraduate and graduate students, taught scientific Python at CUHK, and conducted gravitational-wave data-analysis workshops and seminars at other institutions. I would contribute this experience to CCA's mentoring and software activities while continuing my IGWN software commitments.

Thank you for considering my application.

Sincerely,\\[4pt]
Phurailatpam Hemantakumar
\end{document}
